\documentclass[a4paper,10pt]{article}
\usepackage[utf8]{inputenc}
\usepackage[T1]{fontenc}
\usepackage{lmodern}
\usepackage[margin=1.4cm, top=1.2cm, bottom=1.2cm]{geometry}
\usepackage{hyperref}
\usepackage{xcolor}
\usepackage{enumitem}
\usepackage{graphicx}
\usepackage{titlesec}

% Minimal clean styling
\hypersetup{colorlinks=true, urlcolor=blue!70!black, linkcolor=blue!70!black}
\pagestyle{empty}
\setlength{\parindent}{0pt}

% Section formatting
\titleformat{\section}{\large\bfseries\color{black}}{}{0em}{}[\titlerule]
\titlespacing*{\section}{0pt}{6pt}{4pt}

% Colors
\definecolor{accent}{RGB}{20, 50, 120}

% Role command
\newcommand{\role}[4]{%
  \vspace{2pt}\noindent
  \begin{minipage}[t]{0.70\textwidth}\raggedright\textbf{#1} \textbullet\ \textit{#3}\end{minipage}%
  \hfill
  \begin{minipage}[t]{0.28\textwidth}\raggedleft\footnotesize\color{gray}#2 \textbullet\ #4\end{minipage}%
  \vspace{1pt}%
}
\setlist[itemize]{leftmargin=1.2em, itemsep=1.5pt, topsep=2pt, parsep=0pt}

\begin{document}

% ================= HEADER =================
\begin{minipage}[c]{0.80\textwidth}
  {\Huge\bfseries Prabhavit Gupta}\\[3pt]
  {\color{accent}\large Arquitecto de IA / Senior AI Engineer $\vert$ LLMs, Agentic AI, Model Context Protocol (MCP) $\&$ RAG}\\[6pt]
  {\footnotesize +39 351 610 7807 ~$\vert$~
  \href{mailto:prabhavitg@gmail.com}{prabhavitg@gmail.com} ~$\vert$~
  \href{https://www.linkedin.com/in/prabhavit/}{linkedin.com/in/prabhavit} ~$\vert$~
  \href{https://github.com/prab-gupta}{github.com/prab-gupta} ~$\vert$~
  Venecia, Italia / España (Residente Permanente UE $\vert$ Disponibilidad Inmediata $\vert$ Modalidad Híbrida/Remota)}
\end{minipage}\hfill

\vspace{6pt}

% ================= EXECUTIVE SUMMARY =================
\section{Perfil Profesional / Summary}
\begin{itemize}
  \item \textbf{Arquitectura de IA Generativa \& Sistemas Agenticos en Producción:} Ingeniero de software con ~3 años de experiencia comercial en el diseño y despliegue de soluciones de inteligencia artificial en producción. Especializado en arquitecturas basadas en Large Language Models (modelos open-source y frontier models), agentes autónomos, voicebots/chatbots y automatizaciones inteligentes de procesos.
  \item \textbf{Model Context Protocol (MCP) \& Integración de Herramientas:} Experiencia avanzada desarrollando servidores y clientes basados en el estándar Model Context Protocol (MCP), permitiendo a agentes inteligentes consultar bases de datos relacionales, repositorios documentales y APIs REST externas mediante tool calling y esquemas tipificados (Pydantic).
  \item \textbf{Pipelines RAG Avanzados, Serving de Modelos \& Evaluación:} Especialista en Retrieval-Augmented Generation (RAG) de alta velocidad con PostgreSQL pgvector y Qdrant (chunking semántico, búsqueda híbrida, reranking). Experiencia en serving de inferencia (vLLM, quantización, optimización de latencia y throughput) y suites de evaluación automatizada (RAGAS framework).
\end{itemize}

% ================= TECHNICAL SKILLS =================
\section{Competencias Técnicas / Technical Skills}
\begin{itemize}[leftmargin=1.2em, itemsep=2pt]
  \item \textbf{IA Generativa \& LLMs:} Modelos Open-Source (Llama, Mistral), OpenAI GPT-4o, Anthropic Claude, Prompt Engineering Avanzado, Function/Tool Calling, Structured Outputs (Pydantic / JSON Schema), Guardrails de Seguridad.
  \item \textbf{Model Context Protocol (MCP) \& Agentes:} MCP Server \& Client Development, Conectores de Datos Seguros, LangGraph, LangChain, LlamaIndex, Arquitecturas Multi-Agente, Planificación ReAct, Estado y Memoria.
  \item \textbf{Inferencia \& Optimización:} vLLM, Serving de Inferencia, Cuantización de Modelos, Optimización de GPU, Benchmarking de Latencia, Rendimiento y Consumo de Recursos.
  \item \textbf{RAG, Vector Search \& Bases de Datos:} Retrieval-Augmented Generation (RAG), Ingesta de Documentos, PostgreSQL pgvector, Qdrant, ChromaDB, Búsqueda Híbrida (Dense + Sparse BM25), Caching Semántico (Redis).
  \item \textbf{Ingeniería de Software \& DevOps:} Python (FastAPI, AsyncIO, PyTest, Type Hints), TypeScript, Node.js, Docker, Kubernetes/k3s, Git, GitHub Actions CI/CD, Arquitectura de Microservicios.
  \item \textbf{Evaluación \& Observabilidad:} Framework RAGAS, Precisión/Recall/Fidelidad, OpenTelemetry Tracing, Logging Estructurado, Control de Costes y Presupuesto de Tokens.
\end{itemize}

% ================= PROFESSIONAL EXPERIENCE =================
\section{Experiencia Profesional / Work Experience}

\role{Arquitecto de Soluciones de IA \& Consultor Senior}{2024 -- Presente}{Consultoría Independiente de Ingeniería de IA}{Venecia, IT (Remoto Europa / España)}
\begin{itemize}
  \item Diseñé e implementé arquitecturas complejas de agentes autónomos y pipelines de IA generativa para empresas europeas, traduciendo requisitos operativos en soluciones robustas y escalables.
  \item Desarrollé conectores certificados bajo el protocolo Model Context Protocol (MCP), dotando a agentes LLM de acceso auditable y seguro a bases de datos relacionales y servicios SaaS internos.
  \item Construí motores RAG híbridos de alto rendimiento sobre PostgreSQL pgvector y Qdrant, alcanzando tiempos de respuesta sub-segundo con cero alucinaciones sobre documentación técnica.
  \item Implementé bancos de pruebas y evaluación continua con el framework RAGAS, midiendo fidelidad contextual y relevancia de respuestas ante actualizaciones de modelos e inferencia.
  \item Empaqueté microservicios de IA en contenedores Docker y configuré flujos de integración y despliegue continuo (CI/CD) con observabilidad y tracing distribuido.
\end{itemize}

\vspace{3pt}
\role{Full-Stack AI Software Engineer}{2023 -- 2024}{Tech Automation Systems}{Milán / Remoto, IT}
\begin{itemize}
  \item Desarrollé microservicios asíncronos en Python y TypeScript para el procesamiento, clasificación y extracción estructurada de datos a partir de grandes volúmenes de documentos PDF.
  \item Implementé flujos de indexación vectorial y colas de ingestión en segundo plano con Redis y PostgreSQL, garantizando alta disponibilidad y consistencia de datos.
  \item Creé suites de pruebas unitarias e integración con PyTest, manteniendo altos estándares de cobertura y calidad técnica en cada entrega.
  \item Colaboré activamente con responsables de producto y arquitectos técnicos en la definición de especificaciones de sistemas inteligentes.
\end{itemize}

% ================= NOTABLE PROJECTS =================
\section{Proyectos Destacados / Key Projects}
\begin{itemize}
  \item \textbf{Plataforma Multi-Agente con Servicios MCP (Python + LangGraph):} Sistema colaborativo de agentes para clasificación, extracción y redacción de informes con ejecución de herramientas verificadas vía Model Context Protocol.
  \item \textbf{Motor de Recuperación Documental Empresarial (Python + pgvector):} Motor RAG en producción combinando búsqueda semántica vectorial, búsqueda por palabras clave BM25 y guardrails contra inyecciones de prompt.
\end{itemize}

% ================= EDUCATION & LEGAL STATUS =================
\section{Formación, Idiomas \& Estado Legal}
\begin{itemize}
  \item \textbf{Grado en Filosofía y Economía} -- Universidad Ca' Foscari de Venecia (2018 -- 2022) \\
  \textit{Formación analítica rigurosa en lógica formal, teoría de decisiones y modelado de sistemas complejos.}
  \item \textbf{Idiomas:} Inglés (C1 Profesional, fluido hablado y escrito), Italiano (Competencia profesional), Español (Comprensión y lectura técnica).
  \item \textbf{Autorización Legal de Trabajo:} Residente Permanente de la Unión Europea (Permesso di Soggiorno UE a Tempo Indeterminato, Italia) -- Pleno derecho legal a trabajar en España y en toda la UE sin necesidad de visado o patrocinio.
\end{itemize}

\end{document}
